
\label{app:firefly}

The source code and documentation for this framework are available at
\url{https://github.com/GHeier/FFirefly}.

\section{Overview}

FFirefly is an open-source computational physics framework for condensed matter
research, with emphasis on superconductivity and correlated-electron systems.
The name reflects the project's central design principle: \emph{foreign function
interfaces} (FFI).  A compiled C++ core library exposes a stable \texttt{extern "C"}
interface that both Python (via \texttt{ctypes}) and Julia (via \texttt{ccall}) can
call directly, without language-specific bindings generated by a code-generation
tool.  The result is that all three languages share the same underlying objects and
can pass data between one another across module boundaries.  Each solver project is
therefore free to be written in whichever language is most natural for its
algorithm---tight inner loops in C++, symbolic or array-oriented code in Julia,
ecosystem-heavy calculations (TRIQS, SciPy) in Python---while still consuming and
producing data in a format every other solver understands.

FFirefly is built around a single executable (\texttt{fly.x}) that reads a
configuration file and dispatches to whichever calculation and solver the user
specifies.  Individual solvers are self-contained projects that can be added to or
removed from the build independently; the executable discovers and routes to them
without modification to the surrounding code.

\subsection{Parallelism}

Each solver project may exploit two levels of parallelism.  Within a single node,
OpenMP threads are used to parallelise loops over $\mathbf{k}$-points, frequency
points, and matrix operations.  Across nodes, MPI distributes independent
work units---for example, different $\mathbf{q}$-point slices of a susceptibility
calculation---so that large meshes scale to cluster environments.  Because each
solver is independent, multiple solvers at the same pipeline stage can be launched
simultaneously on separate MPI partitions, then have their outputs consumed by
downstream stages once all are complete.

\subsection{Unified Field Objects and Data Flow}

A central design principle is that every solver reads and writes data through the
same set of typed \emph{field} objects.  A field encapsulates a physical quantity
sampled on a $\mathbf{k}$-mesh and/or a frequency grid, together with the lattice
geometry and tensor structure needed to interpret the values.  Fields serialise to
and from HDF5 files using a fixed schema, and the in-memory representation provides
multilinear interpolation so that any solver can query any field at arbitrary
$(\mathbf{k},\omega)$ points without knowing how the field was produced.

Concretely, the type hierarchy is:
\begin{itemize}
    \item \texttt{Field\_R}, \texttt{Field\_C} --- real and complex scalar fields
        $f(\mathbf{k},\omega) \in \mathbb{R}$ or $\mathbb{C}$.
    \item \texttt{Field\_RM}, \texttt{Field\_CM} --- real and complex matrix fields
        $M_{ij}(\mathbf{k},\omega)$, used for multi-band Hamiltonians, self-energies,
        and gap functions.
    \item \texttt{Field} --- polymorphic wrapper that dispatches to the appropriate
        typed variant at runtime.
\end{itemize}
Because all physics objects---the band dispersion, the susceptibility, the
self-energy, the pairing vertex, the gap function---are instances of the same
interface, the output of any solver is immediately usable as input to any other.
This eliminates format-conversion glue code and makes it straightforward to swap
one method for another at any stage of a pipeline.

\section{Throughput for Computing the Superconducting Gap}
\label{sec:pipeline}

The calculation of the superconducting gap function $\Delta(\mathbf{k},\omega)$
requires a sequence of intermediate quantities, each building on the last.
The following paragraphs describe
each stage and the information it passes forward.

\paragraph{1. Band structure.}
The starting point is the single-particle dispersion $\varepsilon_n(\mathbf{k})$,
obtained by diagonalising the Hamiltonian $H(\mathbf{k})$ at each $\mathbf{k}$-point.
This stage outputs a \texttt{Field\_RM} (the Hamiltonian matrix on the mesh) and
implicitly defines the non-interacting Green's function
$G_0(\mathbf{k},i\omega_n) = [i\omega_n - H(\mathbf{k})]^{-1}$.
The Fermi surface---the iso-energy contour $\varepsilon_n(\mathbf{k})=\mu$---is
extracted at this stage as a list of $\mathbf{k}$-points with associated area
elements, which provide the DOS weights $w_k$ used in downstream surface integrals.

\paragraph{2. Vertex (pairing interaction).}
Given $G_0$, the bare susceptibility
$\chi_0(\mathbf{q},i\nu_m) = -T\sum_{\mathbf{k},n} G_0(\mathbf{k}+\mathbf{q},i\omega_n+i\nu_m)
G_0(\mathbf{k},i\omega_n)$
is computed first and stored as a \texttt{Field\_C}.  The irreducible pairing vertex
is then assembled within the fluctuation-exchange (FLEX) approximation:
\[
    V^{\mathrm{FLEX}}(\mathbf{q},i\nu) =
    U^2 \left[\frac{\chi_0(\mathbf{q},i\nu)}{1 - U\chi_0(\mathbf{q},i\nu)}
    - \tfrac{1}{2}\chi_0(\mathbf{q},i\nu)\right].
\]
The vertex is saved as a \texttt{Field\_C} (or \texttt{Field\_CM} for multi-band
systems) and consumed by both the self-energy and gap-equation stages.

\paragraph{3. Self-energy and renormalization.}
The momentum- and frequency-dependent self-energy
$\Sigma(\mathbf{k},i\omega_n)$ is computed from $V$ and $G_0$ and saved as a
\texttt{Field\_C}.  From $\Sigma$ one extracts the quasiparticle renormalization
weight $Z_n(\mathbf{k}) = [1 - \partial\,\mathrm{Im}\,\Sigma/\partial\omega]^{-1}$,
which enters the gap equation as a mass-renormalization factor and suppresses
the effective pairing strength relative to the bare vertex.  The dimensionless
coupling $\lambda$ is estimated from a Fermi-surface average:
\[
    \lambda = \frac{1}{N_F}\sum_{k,k'} w_k\,Z\,V(\mathbf{k}-\mathbf{k}')\,w_{k'},
\]
providing a fast diagnostic before the full gap equation is solved.

\paragraph{4. Gap equation.}
With $H(\mathbf{k})$, $V(\mathbf{q},i\nu)$, $\Sigma(\mathbf{k},i\omega)$, and $Z$
all available as field objects, the gap equation can be assembled.  Two formulations
are supported:

\begin{itemize}
    \item \textbf{Linearized BCS}: the static limit $V(\mathbf{q})$ and the BCS
        kernel $K(\varepsilon,T) = \tanh(\varepsilon/2T)/2\varepsilon$ define a
        matrix eigenvalue problem on the Fermi surface or over the full Brillouin
        zone.  $T_c$ is identified when the leading eigenvalue reaches unity.
    \item \textbf{Eliashberg}: the full frequency-dependent equation
        retains $V(\mathbf{q},i\nu_m)$, $Z(\mathbf{k},i\omega_n)$, and the
        phonon/electronic renormalization self-consistently.  This is solved as a
        linear operator eigenvalue problem using the Lanczos algorithm, returning
        the gap function $\Delta(\mathbf{k},i\omega_n)$ for each competing pairing
        channel.
\end{itemize}

Because every intermediate result is a field object stored in HDF5, the pipeline
stages are decoupled: any stage can be re-run with a different solver or finer mesh
without invalidating cached results from other stages, and stages whose inputs are
already available can be skipped entirely.

\section{Physics Modules and Algorithmic Dependencies}

\subsection{Electronic Structure (\texttt{hamiltonian})}

\paragraph{Band structure.}
Dispersion relations $\varepsilon_n(\mathbf{k})$ are supplied either via
tight-binding models (hopping parameters $t_0,\ldots,t_{10}$) or via a preloaded
Hamiltonian matrix $H(\mathbf{k})$ that is diagonalised on-the-fly using LAPACK
(\texttt{ssyevd} for real symmetric, \texttt{cheevd} for Hermitian matrices).
Fermi velocities $\mathbf{v}_n(\mathbf{k}) = \nabla_\mathbf{k}\varepsilon_n$ are
computed by finite difference.

\paragraph{Density of States.}
Two methods are available:

\begin{itemize}
    \item \textbf{Gaussian broadening} (\texttt{DOS/gaussian}): evaluates
        $N(\omega)=\sum_{n,\mathbf{k}} g(\omega - \varepsilon_n(\mathbf{k}))$ on a
        uniform $\mathbf{k}$-mesh with a Gaussian $g$ of adjustable width.
    \item \textbf{Tetrahedron method} (\texttt{DOS/tetrahedra}): partitions the
        Brillouin zone into tetrahedra and integrates linearly interpolated bands
        analytically within each tetrahedron, giving spectral weight without
        broadening artefacts at sharp features.
\end{itemize}

\paragraph{Fermi surface.}
The Fermi surface is constructed by the tetrahedron method (\texttt{FS/tetrahedra}),
which locates the iso-energy surface $\varepsilon_n(\mathbf{k})=\mu$ and returns
a list of $\mathbf{k}$-points together with the area element of each face, suitable
for numerical surface integrals of the form
\[
    \langle f \rangle_{\mathrm{FS}} = \int_{\mathrm{FS}}
    \frac{f(\mathbf{k})\,dS_\mathbf{k}}{|\mathbf{v}_F(\mathbf{k})|\,(2\pi)^d}.
\]

\subsection{Many-Body Physics (\texttt{many\_body})}

\paragraph{Self-energy from the vertex (sparse IR, Julia).}
\label{sec:se_sparse_ir}
The self-energy $\Sigma(\mathbf{k}, i\omega_n)$ is computed from the irreducible
vertex $V(\mathbf{q}, i\nu_m)$ and the non-interacting Green's function
$G_0(\mathbf{k}, i\omega_n)$ via the convolution
\[
    \Sigma(\mathbf{k}, i\omega_n) =
    -\frac{1}{N_\mathbf{k}}\sum_{\mathbf{q}} V(\mathbf{q},i\nu_m)
    G_0(\mathbf{k}-\mathbf{q}, i\omega_n - i\nu_m).
\]
To avoid an $\mathcal{O}(N_\omega^2 N_k^2)$ direct double convolution, the
calculation proceeds in three stages:
\begin{enumerate}
    \item Expand $G_0$ and $V$ in the \emph{sparse intermediate representation}
        (IR) basis to obtain a compact set of IR coefficients.
    \item Transform to imaginary time via the IR-to-$\tau$ matrix, then to real
        space via fast Fourier transform (FFTW3), giving $G_0(\mathbf{r},\tau)$
        and $V(\mathbf{r},\tau)$.
    \item Multiply pointwise in $(\mathbf{r},\tau)$ space and transform back to
        $(\mathbf{k}, i\omega_n)$ using inverse FFT and the IR basis.
\end{enumerate}
Matsubara frequencies are sampled on the sparse IR sampling grid, reducing the
number of frequency points from $\mathcal{O}(\beta/\delta\tau)$ to
$\mathcal{O}(\log(\beta\omega_\mathrm{max}))$.

\paragraph{Self-energy via TRIQS / DMFT (Python).}
An alternative route interfaces with the TRIQS library to perform self-consistent
DMFT or IPT calculations at finite temperature.  Green's functions are represented
on discrete Lehmann representation (DLR) Matsubara frequency meshes.  A combined
Bubble+DMFT solver (\texttt{Bubble\_DMFTSolver}) interpolates between the
purely local DMFT self-energy and the momentum-dependent bubble self-energy to
account for non-local correlations.

\paragraph{Vertex and response functions.}
The irreducible pairing vertex $V(\mathbf{q})$ is assembled from the
fluctuation-exchange (FLEX) approximation:
\[
    V^{\mathrm{FLEX}}(\mathbf{q},i\nu) =
    U^2 \left[\frac{\chi(\mathbf{q},i\nu)}{1 - U\chi(\mathbf{q},i\nu)}
    - \frac{1}{2}\chi(\mathbf{q},i\nu)\right],
\]
where $\chi(\mathbf{q},i\nu) = -G_0 * G_0$ is the bare susceptibility computed by
the same sparse IR convolution pipeline described above.

\paragraph{Fermi-surface renormalization.}
The quasiparticle renormalization weight $Z$ and the dimensionless pairing coupling
$\lambda$ are computed on the Fermi surface approximation
(\texttt{renormalization/FS\_approx}):
\[
    \lambda = \frac{\sum_{k,k'} w_k\, U^2 Z\,\chi(\mathbf{k}-\mathbf{k}')\, w_{k'}}
                   {N_F},
\]
where $w_k = dS_k / |\mathbf{v}_F(\mathbf{k})|(2\pi)^d$ are Fermi surface DOS
weights.  An analytic renormalization (\texttt{renormalization/analytic}) and a
direct route from the full self-energy (\texttt{renormalization/from\_sigma}) are
also available.

\subsection{Superconductivity (\texttt{superconductor})}

\paragraph{BCS gap equation (matrix solver, C++).}
The static, linearized BCS gap equation is cast as an eigenvalue problem:
\[
    \Delta_n(\mathbf{k}) =
    -\sum_{\mathbf{k}'} V_{nn'}(\mathbf{k},\mathbf{k}')
    K(\varepsilon_{n'}(\mathbf{k}'),T)\,\Delta_{n'}(\mathbf{k}'),
\]
where $K$ is the BCS kernel $\tanh(\varepsilon/2T)/2\varepsilon$.  The matrix is
assembled on the Fermi surface or over the full Brillouin zone and diagonalised to
find $T_c$ from the leading eigenvalue.  Two methods are available: direct power
iteration for the dominant eigenmode and a full matrix diagonalisation for multi-band
systems.

\paragraph{Eliashberg equation (convolution, Python).}
The full frequency-dependent Eliashberg equation is solved in imaginary-frequency
space:
\[
    Z(k,i\omega_n)\Delta(k,i\omega_n) =
    -\frac{\pi T}{N_k}\sum_{k',m} V(k-k',i\omega_n - i\omega_m)
    \frac{\Delta(k',i\omega_m)}{\sqrt{\omega_m^2 + |\Delta(k',i\omega_m)|^2}},
\]
with a matrix-valued generalisation for multi-band systems.  The Eliashberg kernel
is treated as a \texttt{LinearOperator} and the leading eigenvalues are found with
the Lanczos algorithm (\texttt{scipy.sparse.linalg.eigsh}), which returns the top
eigenvalues without forming the full matrix.  A superconducting instability is
indicated when the leading eigenvalue $\lambda_{\max} \geq 1$.  Green's functions
and vertices are stored on DLR meshes (via TRIQS) and convolutions are performed
with the sparse IR machinery.

\paragraph{Eliashberg equation (hierarchical matrix, Julia).}
A second Eliashberg solver (\texttt{eliashberg/hmatrix}) uses a hierarchical matrix
($\mathcal{H}$-matrix) compression of the kernel to reduce the memory and
computational cost for large $\mathbf{k}$-meshes.

\paragraph{Frequency-dependent BCS with hierarchical matrices (\texttt{bcs\_w/hmatrix}, Julia).}
This solver (\texttt{bcs\_w}) occupies a middle ground between the static BCS
eigenvalue problem and the full self-consistent Eliashberg equation.  It retains
frequency dependence through a perturbative expansion rather than iterating to
self-consistency, which makes it substantially cheaper while still capturing
retardation effects.

The calculation proceeds in two stages.  First, the static pairing kernel
\[
    K^{(0)}_{kk'} = -\sqrt{w_k}\,V(\mathbf{k}-\mathbf{k}')\,\sqrt{w_{k'}}
\]
is compressed into an $\mathcal{H}$-matrix using adaptive cross approximation
(\texttt{HMatrices.jl}).  $\mathcal{H}$-matrix compression exploits the smooth
momentum-space structure of $V$ to approximate off-diagonal blocks by low-rank
matrices, reducing storage and matrix-vector products from $\mathcal{O}(N_k^2)$ to
$\mathcal{O}(N_k \log N_k)$.  The dominant eigenvalues $\lambda_0$ and eigenvectors
of $K^{(0)}$ are then found by Arnoldi/Lanczos iteration (\texttt{KrylovKit.jl}).

Second, a frequency-dependent correction kernel
\[
    K^{(1)}_{kk'} = -\sqrt{w_k}
    \left[\sum_\nu \frac{V(\mathbf{k}-\mathbf{k}',\nu) - V(\mathbf{k}-\mathbf{k}',0)}{\nu}\right]
    \sqrt{w_{k'}}
\]
is compressed into a second $\mathcal{H}$-matrix and evaluated perturbatively.
The first-order correction $\lambda_1 = \langle v_0 | K^{(1)} | v_0 \rangle$ is
combined with the static eigenvalue $\lambda_0$, the quasiparticle weight $Z$, and
a Coulomb pseudopotential $\mu^*$ to give an effective coupling
\[
    \lambda_\mathrm{eff} = \frac{\lambda_0\,(1 - \mu^*)}{1/Z + \lambda_1},
\]
from which $T_c$ is estimated via the McMillan-like expression
$T_c = 1.134\,\omega_c\,\exp(-1/\lambda_\mathrm{eff})$.
This formulation connects directly to the renormalization quantities computed in
the \texttt{many\_body} module: $Z$ is read from the self-energy field object and
$\lambda_1$ encodes the same retardation integral that the Eliashberg equation
treats self-consistently.

\bigskip

